% Generated draft: requires mathematical and source review.
\documentclass[11pt]{article}
\usepackage[margin=25mm]{geometry}
\usepackage{fontspec}
\setmainfont{Latin Modern Roman}
\usepackage{amsmath,amssymb,mathtools,unicode-math}
\setmathfont{Latin Modern Math}
\usepackage{enumitem,needspace,xurl,hyperref}
\setlength{\parindent}{0pt}
\setlength{\parskip}{5pt}
\setlength{\emergencystretch}{3em}
\newcommand{\N}{\mathbb{N}}
\newcommand{\Z}{\mathbb{Z}}
\newcommand{\Q}{\mathbb{Q}}
\newcommand{\R}{\mathbb{R}}
\newcommand{\C}{\mathbb{C}}
\newcommand{\F}{\mathbb{F}}
\newcommand{\A}{\mathbb{A}}
\newcommand{\PP}{\mathbb{P}}
\newcommand{\E}{\mathbb{E}}
\newcommand{\eps}{\varepsilon}
\newcommand{\dd}{\,\mathrm{d}}
\begin{document}
\section*{UnsolvedMath ID: 30002265}
\textbf{Source problem code:} catalogue entry
\subsection{Definitions and mathematical statement}
Interpret the dimension as a positive integer \(d\), and interpret the \(n\) points as distinct. These conventions are explicit here because the supplied record does not separately specify them. For integers \(d\geq 1\) and \(n\geq d+1\), let
\[
\mathcal S_{d,n}=\left\{S\subseteq\R^d: |S|=n,\ \max_{x,y\in S}\|x-y\|_2=1\right\}.
\]
For \(S\in\mathcal S_{d,n}\), its diameter graph \(D(S)\) has vertex set \(S\) and edge set
\[
E(D(S))=\left\{\{x,y\}\subseteq S:x\ne y,\ \|x-y\|_2=1\right\}.
\]
A \((d-1)\)-dimensional simplex in this graph is interpreted as an unordered set of \(d\) vertices forming a complete subgraph. For \(d\geq2\), these vertices form a regular Euclidean simplex of dimension \(d-1\); for \(d=1\), each singleton is a \(0\)-dimensional simplex. Define
\[
s_{d-1}(S)=\#\left\{T\subseteq S:|T|=d,\ \|x-y\|_2=1\text{ for all distinct }x,y\in T\right\}.
\]
The supplied corrected question, attributed in the record to Schur's diameter-simplex conjecture, is whether
\[
\forall d\in\Z_{\geq1}\setminus\{3\}\quad
\forall n\in\Z_{\geq d+1},\qquad
\max_{S\in\mathcal S_{d,n}}s_{d-1}(S)=n.
\]
Thus the proposed equality includes both an upper bound of \(n\) for every admissible configuration and attainment by some admissible configuration for each parameter pair.

The original supplied formulation does not exclude \(d=3\). The corrected formulation excludes that dimension because the supplied dated curation reports it as proved there. The same curation describes other dimensions as open and mentions conditional results without stating their hypotheses or conclusions. These are supplied source qualifications, not independently verified mathematical or literature assessments; in particular, the claimed unresolved range needs review of the intended lower bound on \(d\).
\subsection{Short English statement}
For each positive integer dimension \(d\ne3\) and each \(n\geq d+1\), consider \(n\) distinct Euclidean points of diameter exactly \(1\). Is the largest possible number of complete \(d\)-vertex subgraphs in their diameter graph exactly \(n\)? The dimension convention, the original inclusion of \(d=3\), and the supplied dated status assessment require human review.
\subsection{Sources}
\begin{itemize}
\item Supplied local UnsolvedMath registry record, permanent ID \(30002265\), title: Maximum Simplices in Euclidean Diameter Graphs; exact source problem code: \texttt{catalogue entry}. The record belongs to the Oberwolfach Reports Open Problems collection and gives proposed year \(2013\); no individual proposer is supplied.
\item Supplied report citation: \emph{Graph Theory} (2013), Oberwolfach Reports, pp.\ 67--128, DOI \texttt{10.4171/owr/2013/02}. Supplied source URL and literature-source link, labelled Graph Theory, Oberwolfach Report 02/2013: \url{https://doi.org/10.4171/owr/2013/02}. The supplied background describes Oberwolfach Reports as documenting mathematical workshops and their open-problem discussions.
\item Attribution and status: the supplied record identifies the question as Schur's diameter-simplex conjecture and labels it partially solved. Its literature assessment reports a proof in dimension three, an open status in other dimensions, and conditional results in the report. Literature checked at: \texttt{2026-08-21}. This is supplied dated curation, not independently verified; the record itself requests independent verification of time-sensitive status and cited claims.
\item The supplied statement-verification field says the official report supports the corrected formulation and carries the label \texttt{corrected\_verified}. Those verification assertions are reproduced as registry provenance, not adopted as independent verification in this draft.
\end{itemize}
\end{document}
